% !TeX root = ../../main.tex
% sections/app/appendix_cost_minimization.tex

\chapter{費用最小化問題からの需要関数と価格指数の導出（第1段階）}
\label{chap:appendix_cost_minimization}

本付録の目的は、総名目生産額が、集計産出量と真の価格指数の単純な積で表せること、そしてその関係を担保するために集計産出量が必然的にCES集計の形でなければならないことを証明することである。

\section*{定理2：個別財への需要関数と真の価格指数}
\label{sec:appendix_cost_minimization_theorem2}

家計\(h\)が、ある時点\(t\)において、所与の量の国内財バスケット\(c_t^{h\to H}=\left[\sum_{h'\in H}(c_t^{h\to h'})^{\frac{\theta^H-1}{\theta^H}}\right]^{\frac{\theta^H}{\theta^H-1}}\)を、個別財の価格\(\{p_t^{h'}\}\)の下で費用最小化的に購入するとき、以下の関係が成立する。
\begin{enumerate}
    \item\label{itm:appendix_cost_minimization_theorem2_1}家計\(h\)の個別財\(h'\)への需要関数は、次式で与えられる。
    \begin{equation}
    c_t^{h\to h'}=\left(\frac{p_t^{h'}}{p_t^H}\right)^{-\theta^H}c_t^{h\to H}
    \label{eq:appendix_cost_minimization_theorem2_individual_demand}
    \end{equation}
    \item\label{itm:appendix_cost_minimization_theorem2_2}上記の需要関数に現れる\(p_t^H\)は、国内財バスケットの「真の価格指数」であり、個々の財価格から以下のように定義される。この値は全ての国内家計で共通である。
    \begin{equation}
    p_t^H=\left[\sum_{h'\in H}(p_t^{h'})^{1-\theta^H}\right]^{\frac{1}{1-\theta^H}}
    \label{eq:appendix_cost_minimization_theorem2_price_index}
    \end{equation}
\end{enumerate}

\section*{証明}
\label{sec:appendix_cost_minimization_proof}

この定理を、自国家計\(h\)の場合について証明する。外国の家計\(f\)についても全く対称的な手順で証明可能である。

\subsection*{1.費用最小化問題の設定}
\label{sec:appendix_cost_minimization_sub1}
家計が直面する問題は、時点\(t\)において、所与の国内財バスケット\(c_t^{h\to H}\)を達成するための総費用を最小化することである。
\begin{itemize}
    \item\label{itm:appendix_cost_minimization_sub1_obj} \textbf{最小化対象}:
    \begin{equation}
    \min_{\{c_t^{h\to h'}\}_{h'\in H}}\sum_{h'\in H}p_t^{h'}c_t^{h\to h'}
    \label{eq:appendix_cost_minimization_sub1_objective}
    \end{equation}
    \item\label{itm:appendix_cost_minimization_sub1_const} \textbf{制約条件}:
    \begin{equation}
    c_t^{h\to H}=\left[\sum_{h'\in H}(c_t^{h\to h'})^{\frac{\theta^H-1}{\theta^H}}\right]^{\frac{\theta^H}{\theta^H-1}}
    \label{eq:appendix_cost_minimization_sub1_constraint}
    \end{equation}
\end{itemize}

\subsection*{2.ラグランジュ関数と一階の条件（FOC）}
\label{sec:appendix_cost_minimization_sub2}
この問題を解くために、ラグランジュ関数\(\mathcal{L}^{h\to H}\)を設定する。ラグランジュ乗数を\(\mu_t^{h\to H}\)とする。
\begin{equation}
\mathcal{L}^{h\to H}=\sum_{h'\in H}p_t^{h'}c_t^{h\to h'}-\mu_t^{h\to H}\left(\left[\sum_{h'\in H}(c_t^{h\to h'})^{\frac{\theta^H-1}{\theta^H}}\right]^{\frac{\theta^H}{\theta^H-1}}-c_t^{h\to H}\right)
\label{eq:appendix_cost_minimization_sub2_lagrangian}
\end{equation}
このラグランジュ関数を、任意の個別財\(c_t^{h\to h'}\)で偏微分し、ゼロと置くことで一階の条件（FOC）が得られる。
\begin{align}
\frac{\partial\mathcal{L}^{h\to H}}{\partial c_t^{h\to h'}}=p_t^{h'}-\mu_t^{h\to H}\cdot\frac{\theta^H}{\theta^H-1}\left[\sum_{i\in H}(c_t^{h\to i})^{\frac{\theta^H-1}{\theta^H}}\right]^{\frac{\theta^H}{\theta^H-1}-1}\cdot\frac{\theta^H-1}{\theta^H}(c_t^{h\to h'})^{\frac{\theta^H-1}{\theta^H}-1}&=0 \label{eq:appendix_cost_minimization_sub2_foc_step1} \\
p_t^{h'}&=\mu_t^{h\to H}\cdot(c_t^{h\to H})^{\frac{1}{\theta^H}}\cdot(c_t^{h\to h'})^{-\frac{1}{\theta^H}} \label{eq:appendix_cost_minimization_sub2_foc_step2} \\
p_t^{h'}&=\mu_t^{h\to H}\left(\frac{c_t^{h\to H}}{c_t^{h\to h'}}\right)^{\frac{1}{\theta^H}}
\label{eq:appendix_cost_minimization_sub2_foc_final}
\end{align}

\subsection*{3.FOCからの需要関数の導出}
\label{sec:appendix_cost_minimization_sub3}
上記で得られたFOCを、個別財への需要量\(c_t^{h\to h'}\)について解く。
\begin{align}
    \frac{p_t^{h'}}{\mu_t^{h\to H}}&=\left(\frac{c_t^{h\to H}}{c_t^{h\to h'}}\right)^{\frac{1}{\theta^H}} \label{eq:appendix_cost_minimization_sub3_demand_step1} \\
    \left(\frac{p_t^{h'}}{\mu_t^{h\to H}}\right)^{\theta^H}&=\frac{c_t^{h\to H}}{c_t^{h\to h'}}
    \label{eq:appendix_cost_minimization_sub3_demand_step2}
\end{align}
これによりラグランジュ乗数\(\mu_t^{h\to H}\)を含む需要関数が導出される。
\begin{equation}
c_t^{h\to h'}=\left(\frac{p_t^{h'}}{\mu_t^{h\to H}}\right)^{-\theta^H}c_t^{h\to H}
\tag{\ref{eq:model_households_indices_aggregation_stage1a_home_c_h_H_problem_c_h_h_eq}}
\end{equation}

\subsection*{4.真の価格指数\(p_t^H\)の導出}
\label{sec:appendix_cost_minimization_sub4}
次に、ラグランジュ乗数\(\mu_t^{h\to H}\)の具体的な形を導出する。ステップ3で得た需要関数を、制約条件であるCES集計式に代入する。
\begin{equation}
c_t^{h\to H}=\left[\sum_{h'\in H}\left(\left(\frac{p_t^{h'}}{\mu_t^{h\to H}}\right)^{-\theta^H}c_t^{h\to H}\right)^{\frac{\theta^H-1}{\theta^H}}\right]^{\frac{\theta^H}{\theta^H-1}}
\label{eq:appendix_cost_minimization_sub4_ces_substitution}
\end{equation}
式を整理していく。
\begin{align}
    c_t^{h\to H}&=\left[\sum_{h'\in H}\left(\frac{p_t^{h'}}{\mu_t^{h\to H}}\right)^{1-\theta^H}(c_t^{h\to H})^{\frac{\theta^H-1}{\theta^H}}\right]^{\frac{\theta^H}{\theta^H-1}} \label{eq:appendix_cost_minimization_sub4_ces_step1} \\
    &=\left[(c_t^{h\to H})^{\frac{\theta^H-1}{\theta^H}}(\mu_t^{h\to H})^{\theta^H-1}\sum_{h'\in H}(p_t^{h'})^{1-\theta^H}\right]^{\frac{\theta^H}{\theta^H-1}} \label{eq:appendix_cost_minimization_sub4_ces_step2} \\
    &=(c_t^{h\to H})\cdot(\mu_t^{h\to H})^{\theta^H}\cdot\left[\sum_{h'\in H}(p_t^{h'})^{1-\theta^H}\right]^{\frac{\theta^H}{\theta^H-1}}
    \label{eq:appendix_cost_minimization_sub4_ces_final}
\end{align}
両辺の\(c_t^{h\to H}\)を消去し、\(\mu_t^{h\to H}\)について解く。
\begin{align}
    1&=(\mu_t^{h\to H})^{\theta^H}\left[\sum_{h'\in H}(p_t^{h'})^{1-\theta^H}\right]^{\frac{\theta^H}{\theta^H-1}} \label{eq:appendix_cost_minimization_sub4_mult_step1} \\
    (\mu_t^{h\to H})^{-\theta^H}&=\left[\sum_{h'\in H}(p_t^{h'})^{1-\theta^H}\right]^{\frac{\theta^H}{\theta^H-1}} \label{eq:appendix_cost_minimization_sub4_mult_step2} \\
    \mu_t^{h\to H}&=\left[\sum_{h'\in H}(p_t^{h'})^{1-\theta^H}\right]^{\frac{1}{1-\theta^H}}
    \label{eq:appendix_cost_minimization_sub4_mult_final}
\end{align}
このラグランジュ乗数\(\mu_t^{h\to H}\)は、個々の家計\(h\)に依存しない共通の価格リストのみで決定されるため、全ての国内家計で共通の値をとる。
本稿では、この家計が直面する真の価格指数を\(p_t^H\)と定義する。
\begin{equation}
p_t^H\equiv\mu_t^{h\to H}=\left[\sum_{h'\in H}(p_t^{h'})^{1-\theta^H}\right]^{\frac{1}{1-\theta^H}}
\label{eq:appendix_cost_minimization_sub4_price_index_def}
\end{equation}
この\(p_t^H\)をステップ3の需要関数に代入することで、定理の項目1が証明される（証明終）。